\chapter{ML in the master's curricula}\label{ch-02}

This slide deck is an ``interlude'' of the first lecture, which the
professor asks students to read on their own. It describes how the
Machine Learning course fits into the \textbf{Master's Degree in
Computer Science} at the University of Pisa and its curricula. It
contains no exam material, but it helps to understand the role of the
course within the study programme.

\section{The master's
curricula}\label{i-curricula-della-magistrale}

Since 2017 the master's degree in Computer Science has been organised
into four curricula:

\begin{itemize}
\tightlist
\item
  \textbf{Artificial Intelligence} (AI);
\item
  \textbf{Data and Knowledge}, since 2021 \textbf{Big Data
  Technologies} (BD);
\item
  \textbf{ICT Solutions Architect} (ICT);
\item
  \textbf{Software: Programming, Principles, and Technologies} (SW).
\end{itemize}

The organisation into curricula allows students to specialise in an
area (with a diploma supplement certifying the curriculum), offers
basic methodological courses at the beginning of the programme and
aims to show how two more years of study are useful for one's
professional future. Regulations, exam instructions and procedures for
changing curriculum or study plan are on the department website.

\section{The role of Machine
Learning}\label{il-ruolo-del-machine-learning}

ML is connected to \textbf{all four} curricula, but it is a
\textbf{core course of the AI curriculum}, where it is part of the
methodological basis for building adaptive and intelligent systems,
together with:

\begin{itemize}
\tightlist
\item
  \emph{Artificial Intelligence Fundamentals} (6 ECTS);
\item
  \emph{Computational Mathematics for Learning and Data Analysis} (9
  ECTS);
\item
  \emph{Machine Learning} (9 ECTS);
\item
  \emph{Parallel and Distributed Systems: Paradigms and Models} (9
  ECTS).
\end{itemize}

The other core courses of the AI curriculum are \emph{Human Language
Technologies}, \emph{Intelligent Systems for Pattern Recognition} and
\emph{Smart Applications}, alongside groups of elective exams
(Algorithm Engineering, Data Mining, Mobile and Cyber-Physical Systems,
Information Retrieval, Computational Neuroscience, Robotics, Semantic
Web, etc.).

ML is the pivot of the \textbf{intelligent systems area}: many master's
courses (Human Language Technologies, Intelligent Systems for Pattern
Recognition, Smart Applications, Computational Neuroscience, Robotics,
Data Mining, Information Retrieval, Bioinformatics, Big Data
Analytics\ldots) directly use its methods.

\begin{boxnote}{Study plan of the AI curriculum (since 2021)}

60 ECTS of core courses, 27 ECTS from groups of elective exams (one
9-ECTS exam and three 6-ECTS exams) and at least 9 ECTS of free choice,
which can also be covered with two 6-ECTS exams.

\end{boxnote}

\section{Synergy with Computational
Mathematics}\label{sinergia-con-computational-mathematics}

The professor recommends following \textbf{ML} and \textbf{CM}
(\emph{Computational Mathematics for Learning and Data Analysis}) in
parallel: the former provides the methods, the latter the mathematical
foundations, and the two courses stimulate each other. CM is however
not a prerequisite: many students from other degree programmes or
curricula have taken ML without CM.

\section{Notes for specific categories of
students}\label{note-per-categorie-particolari-di-studenti}

\begin{boxwarning}{For those coming from the Pisa bachelor's degree (IIA course)}

The first six lectures will seem ``easy'', but they \textbf{do not
coincide} with the contents of IIA: they contain new parts and
different mathematical concepts. The change of perspective is sharp: in
IIA simple models were the \emph{tool} for introducing ML concepts,
here state-of-the-art models and principles are the \emph{object} of
study. From lecture 5--6 onwards everything is new.

\end{boxwarning}

For students of the old \textbf{AA1} exam (320AA, 6 ECTS), which no
longer exists: those who need to switch from AA1 to ML (9 ECTS) take a
supplementary exam on a list of topics indicated by the professor;
those who need to do the opposite must contact the lecturer.

For any information the reference is Prof.~Alessio Micheli
(micheli@di.unipi.it), of the \emph{Computational Intelligence \&
Machine Learning} group of the Department of Computer Science.
